\section{Scope, conventions, and the three meanings of BCH}\label{sec:scope}
The question that organizes this article is elementary to state: given two
quantities $X,Y$ whose products need not commute, can the product $e^Xe^Y$ be
written as a single exponential, and what exactly is its exponent? The central
answer is
\begin{equation}\label{eq:main}
 e^Xe^Y=e^{Z(X,Y)},\qquad
 Z(X,Y)=X+Y+\tfrac12[X,Y]+\tfrac1{12}[X,[X,Y]]+\tfrac1{12}[Y,[Y,X]]+\cdots.
\end{equation}
The substantive assertion is not the existence of a formal logarithm, which is
immediate in a completed free associative algebra. It is that every homogeneous
part of that logarithm is a \emph{Lie polynomial}, a linear combination of
nested commutators, and that its coefficients can be specified completely.

The equation $e^Xe^Y=e^Z$ is not, by itself, a definition of a unique
logarithm. The exponential can fail to be injective, and a product of
exponentials need not be an exponential \emph{in the specified Lie algebra}.
What is universal is the formal series
\begin{equation}\label{eq:formalbch}
 Z(X,Y)=\BCH(X,Y)=\log(e^Xe^Y)=\sum_{n\geq1}Z_n(X,Y),
\end{equation}
where $Z_n$ has total degree $n$. A principal matrix logarithm, a continued
branch of a logarithm, and a convergent BCH series are related but distinct
objects.

For definiteness we state at the outset, in both settings, the assertion that
the rest of the article proves and extends. Here $Z_n(X,Y)$ is the
homogeneous component of total degree $n$ produced by the degree recursion
\eqref{eq:degreerec}, and a \emph{closed Lie subalgebra} of a Banach algebra
is a closed linear subspace closed under the commutator.

\begin{theorem}[Baker--Campbell--Hausdorff]\label{thm:bch}
\begin{enumerate}[label=(\alph*)]
\item \emph{(Formal form, \eqref{eq:formalbch}.)} In the free associative
algebra $\kk\langle X,Y\rangle$ over a field $\kk$ of characteristic zero,
for every $n\geq1$ the element $Z_n$ is homogeneous of degree $n$, is a Lie
polynomial in $X$ and $Y$, and satisfies
\[
 Z_n=\sum_{|w|=n}\coeff{w}Z_n\;w
 =\frac1n\sum_{|w|=n}\coeff{w}Z_n\;R(w),
\]
where $R$ is the right-nested bracket \eqref{eq:rightbracket} and the
coefficients $\coeff{w}Z_n=c(w)$ are given by the finite sum
\eqref{eq:wordblocks}.
\item \emph{(Analytic form, \eqref{eq:main}.)} Let $\cA$ be a unital Banach
algebra over $\kk=\R$ or $\C$ with submultiplicative norm, and let
$X,Y\in\cA$ satisfy $\norm X+\norm Y<\log2$. Then $\sum_{n\geq1}Z_n(X,Y)$
converges absolutely; its sum $Z$ is the Mercator logarithm
$\log(e^Xe^Y)=\sum_{k\geq1}\frac{(-1)^{k-1}}k(e^Xe^Y-I)^k$ and satisfies
$e^Z=e^Xe^Y$; an element $Z_0$
with $\norm{Z_0}<\log2$ satisfies $e^{Z_0}=e^Xe^Y$ if and only if $Z_0=Z$;
every $Z_n(X,Y)$, and therefore $Z$ itself, lies in every closed Lie
subalgebra of $\cA$ that contains $X$ and $Y$; and
\[
 Z_1=X+Y,\qquad Z_2=\tfrac12[X,Y],\qquad
 Z_3=\tfrac1{12}[X,[X,Y]]+\tfrac1{12}[Y,[Y,X]].
\]
\end{enumerate}
\end{theorem}

Part (a) is \cref{thm:wordcut,thm:formalbch,thm:dynkin}, part (b) is
\cref{thm:convergence,prop:uniquelog,prop:closedlie,prop:lowdegree}; the
proofs are given there. Neither part asserts that $Z$ is a principal
logarithm, and neither is claimed outside its stated setting. In
particular part (b) does not assert $\norm Z<\log2$: the bound proved in
\cref{cor:bchlog} is $\norm Z\leq-\log(2-e^{\norm X+\norm Y})$, which
exceeds $\log2$ as $\norm X+\norm Y$ approaches $\log2$, and the
uniqueness clause is stated accordingly.

\subsection{Coverage target and what is not claimed}
The coverage target is the English Wikipedia article
\emph{Baker--Campbell--Hausdorff formula} in the revision with identifier
\texttt{1368909113}, last edited on 11 August 2026 and consulted on
16 September 2026 \cite{wiki}. Every displayed mathematical identity and every
mathematical assertion stated in its main text or its mathematical notes is
addressed in this article; \cref{app:coverage} gives a claim-by-claim
concordance. The article is used as a list of questions, not as a substitute
for proofs. Historical priority statements and application surveys are
bibliographical claims rather than theorems. A theorem merely named in a link,
such as the Stone--von Neumann uniqueness theorem or the Golden--Thompson
inequality, is not thereby stated on the page and is not part of the coverage
target; the actual quantum-mechanical exponential formulas on the page are
proved directly in \cref{sec:quantum}.

Several assertions on the page require correction or explicit hypotheses
rather than proof as written. The coframe matrix in exponential coordinates
is $W=\phi(\ad_X)$ and not the adjoint matrix $\ad_X$, which is always
singular in positive dimension; a Killing form need not be a metric; an
arbitrary pullback can be degenerate; the trace identity concerns the BCH
branch of the logarithm; and a commutator relation on a dense domain does not
by itself imply exponentiated Weyl relations. In each case we prove the valid
statement and exhibit a counterexample to the unqualified one.

The algebraic proofs include the structural results they use, in
particular the ordered Poincar\'e--Birkhoff--Witt theorem. Standard
foundations of analysis are assumed without proof: the inverse function
theorem, existence and uniqueness for ordinary differential equations with
analytic or bounded-linear right-hand sides, the identity theorem for
analytic functions, the contraction mapping theorem, Lebesgue's dominated
and monotone convergence theorems, uniqueness of Fourier coefficients for
continuous functions, and, in \cref{sec:quantum}, the spectral calculus of
self-adjoint operators, Stone's theorem on one-parameter unitary groups,
and the basic criterion for essential self-adjointness. No result about
the BCH formula itself is assumed.

The present article combines three earlier treatments of the same coverage
target that were prepared independently. Where those treatments proved the
same result by different routes, we keep the route that is most complete or
most elementary; where one contained a construction absent from the others,
that construction is included here. All proofs are written out in full.

\subsection{Three distinct settings}
Conflating the following settings is the source of most apparent paradoxes
surrounding the formula.
\begin{description}[leftmargin=1.5em,style=nextline]
\item[Formal algebra.] $X$ and $Y$ have positive degree in a graded algebra
over a field $\kk$ of characteristic zero. Every identity is an identity of
formal series, and every fixed-degree coefficient is a finite sum. No numerical
convergence is intended or asserted.
\item[Local analysis.] $X,Y$ belong to a real or complex unital Banach algebra
$\cA$ with a submultiplicative norm, or to a finite-dimensional Lie algebra,
and are sufficiently small. The series converges and specifies the logarithm
branch issuing from zero. Finite matrices are the principal example.
\item[Global identities under special relations.] Commutation relations may
allow the series to terminate or to be replaced by a closed expression. An
independent differential argument can then prove an exponential identity far
beyond the convergence domain of its Taylor series. Such an identity need not
identify a principal logarithm, and for unbounded operators domains cannot be
suppressed.
\end{description}
A matrix may have logarithms even though its BCH series diverges; a closed
expression may continue the local BCH logarithm without making the original
series convergent; and a logarithm need not belong to the Lie algebra generated
by its two inputs. Each of these phenomena is exhibited by an explicit example
below.

\subsection{Conventions}
Throughout, $[A,B]=AB-BA$ in an associative algebra, $\ad_AB=[A,B]$, and
compositions of adjoint maps are read from right to left. Associativity alone
gives antisymmetry and the Jacobi identity: expanding
$[A,[B,C]]+[B,[C,A]]+[C,[A,B]]$ produces each of the six ordered products
$ABC,ACB,BAC,BCA,CAB,CBA$ twice with opposite signs. Two consequences used
constantly are the derivation rule and its adjoint form,
\begin{equation}\label{eq:derivation}
 [A,[B,C]]=[[A,B],C]+[B,[A,C]],\qquad [\ad_A,\ad_B]=\ad_{[A,B]}.
\end{equation}
For a nonempty word $w=a_1\cdots a_n$ in the letters $X,Y$ its
\emph{right-nested bracket} is
\begin{equation}\label{eq:rightbracket}
 \Rb{a_1\cdots a_n}=[a_1,[a_2,\ldots,[a_{n-1},a_n]\ldots]],
 \qquad \Rb{a}=a,\qquad R(1)=0.
\end{equation}
Thus $\Rb{XXY}=[X,[X,Y]]=\ad_X^2Y$ and $\Rb{YYX}=\ad_Y^2X$. A symbol such as
$X^rY^s$ inside $R(\cdot)$ describes a word with $r$ letters $X$ followed by
$s$ letters $Y$, not a commutator of associative powers; empty runs are
omitted. We write $Z_n$ for the part of $Z$ of total degree $n$, and $Z_{p,q}$
for the part of degree $p$ in $X$ and $q$ in $Y$. For an associative
polynomial $P$ and a word $w$, $\coeff{w}P$ denotes the coefficient of $w$ in
$P$; it is never a commutator. The notation $O_{\deg}(m)$ denotes terms of
total degree at least $m$, and $O(Y^m)$ denotes terms containing at least $m$
occurrences of $Y$, with no claim of uniformity as $X$ varies.

The scalar functions that organize the entire subject are
\begin{equation}\label{eq:scalarfunctions}
\begin{gathered}
 \phi(z)=\frac{1-e^{-z}}{z},\qquad
 \phi(-z)=\frac{e^{z}-1}{z},\qquad
 \beta(z)=\frac{z}{1-e^{-z}},\\
 \gamma(z)=\beta(-z)=\frac{z}{e^z-1},\qquad
 \psi(x)=\frac{x\log x}{x-1}.
\end{gathered}
\end{equation}
The first two are entire, with value $1$ at zero. The functions $\beta$ and
$\gamma$ are meromorphic with a removable singularity at zero and simple poles
at the nonzero points of $2\pi\ii\Z$. The function $\psi$ is used near $x=1$,
with the logarithm branch analytic there and $\psi(1)=1$. An expression such
as $\phi(M)$ or $\beta(\ad_X)$ always denotes functional calculus, formal or
analytic, and never presupposes that $M$ or $\ad_X$ is invertible.

Bernoulli numbers appear in the convention $B_1^+=+\tfrac12$ of the source
page; we write $\bplus{n}=B_n^+/n!$ for the Taylor coefficients of $\beta$ and
$\bminus{n}=(-1)^n\bplus{n}=B_n^-/n!$ for those of $\gamma$.

\subsection{A guide to the results}
Sections \ref{sec:formal}--\ref{sec:dynkin} establish the algebraic theorem
and every coefficient: a finite formula for each associative word coefficient,
Dynkin's commutator formula, a finite formula for each bidegree, and the
complete polynomial through degree six with a finite certificate. Sections
\ref{sec:differential}--\ref{sec:integral} give the differential
constructions: the Campbell identity, the differential of the exponential,
Bernoulli coefficients, the Poincar\'e integral, every order in the number of
occurrences of $Y$, and two total-degree recursions.
\Cref{sec:convergence} proves quantitative convergence estimates, the trace
identity, and the nonconvergence example. \Cref{sec:liegroups} treats Lie
groups, and \cref{sec:special} the exactly summable cases.
\Cref{sec:zassenhaus,sec:splitting} treat the Zassenhaus factorization and the
product formulas. \Cref{sec:geometry,sec:quantum} treat exponential
coordinates and the quantum-mechanical identities. \Cref{sec:computation}
describes the exact computations. The appendices contain the ordered
Poincar\'e--Birkhoff--Witt theorem, coefficient certificates and further
complete degrees, and the coverage concordance.

A reader interested chiefly in coefficients can use the nonrecursive formulas
\eqref{eq:wordcut}, \eqref{eq:dynkin}, \eqref{eq:bidegree},
\eqref{eq:operatorblock}, and \eqref{eq:treeformula} directly.
